\chapter{Hydrological Twin - Technical specifications} \label{chap-tech-surf}

\phantomsection
\fancyhead[LE]{\thepage} 
\fancyhead[RE]{\textbf{Hydrological Twin - Technical Specification}}
\fancyhead[RO]{\thepage} 
\fancyhead[LO]{\textbf{Hydrological Twin - Technical Specification}}
\renewcommand{\headrulewidth}{1pt}

\setlength{\parindent}{0cm} % no indent
\setlength{\parskip}{9pt}

% --------------------------------------------------------------
\section{Goal}

The \texttt{HydrologicalTwin} (HT) is a monolithic backend façade for CaWaQS-ViZ.  
It is designed as the \emph{only} backend entry point used by the QGIS plugin.

HT reuses the existing CaWaQS-ViZ code base without rewriting it, including:
\begin{itemize}
    \item \texttt{Compartment}, \texttt{Mesh}, \texttt{Observations}, \texttt{Extraction}
    \item \texttt{Manage.Temporal}, \texttt{Manage.Budget}, \texttt{Manage.Spatial}
\end{itemize}

It exposes a clean and stable API for:
\begin{itemize}
    \item value extraction from CaWaQS binary outputs,
    \item temporal aggregation,
    \item mass balance and budget analysis,
    \item spatial operators (catchment cells, aquifer outcropping).
\end{itemize}

It also supports persistence via \texttt{.pkl} files for bancarisation.

This architecture deliberately does not rewrite the existing code base; instead, it wraps it and formalizes a clear interface for all front-end clients.

\section{Integration with Existing CaWaQS-ViZ Code}

\subsection{Package Layout}

The new \texttt{ht} subpackage is introduced inside the existing \texttt{cawaqsviz} package:

\begin{verbatim}
cawaqsviz/
    Compartment.py
    Mesh.py
    Observations.py
    Extraction.py
    Manage.py
    interface/
        Config.py
    ht/
        __init__.py
        hydrological_twin.py
        api_types.py
        persistence.py
\end{verbatim}

The existing modules (\texttt{Compartment}, \texttt{Mesh}, \texttt{Observations}, \texttt{Extraction}, \texttt{Manage}) remain unchanged.

\subsection{HydrologicalTwin Construction}

The \texttt{HydrologicalTwin} is built from the same configuration objects already used by \texttt{Compartment}:
\begin{itemize}
    \item \texttt{ConfigGeometry}
    \item \texttt{ConfigProject} (e.g. regime)
    \item \texttt{out\_caw\_directory} (root of CaWaQS outputs)
    \item \texttt{obs\_directory} (root of observation data)
    \item \texttt{temp\_directory} (optional, defaults to CaWaQS outputs directory)
\end{itemize}

The constructor stores these references and instantiates the service layers:
\begin{itemize}
    \item \texttt{Manage.Temporal()}
    \item \texttt{Manage.Budget()}
    \item \texttt{Manage.Spatial()}
\end{itemize}

\section{Domain Model and Encapsulation}

\subsection{Compartments}

HT maintains an internal mapping:
\[
\texttt{self.compartments: Dict[int, Compartment]}
\]

\begin{itemize}
    \item \texttt{build\_compartment(id\_compartment)} calls the existing \texttt{Compartment} constructor.
    \item \texttt{get\_compartment(id\_compartment)} returns an existing instance or builds it on demand.
\end{itemize}

The GUI must never construct \texttt{Compartment} objects directly.  
All access goes through \texttt{HydrologicalTwin.get\_compartment()} or higher-level HT methods.

\subsection{Observations, Mesh, Extraction}

These classes remain internal implementation details of \texttt{Compartment}.  
They are not exposed to the GUI. All access goes through explicit methods on \texttt{HydrologicalTwin}.

\section{HydrologicalTwin Backend API}

The HT API consists of the public methods of the \texttt{HydrologicalTwin} class.  
Each method is a thin wrapper around existing domain logic.

\subsection{\texttt{extract\_values(...)}}

\textbf{Purpose:}  
Read CaWaQS binary outputs, build a time-series DataFrame, and optionally write a CSV.

\textbf{Implementation:}
\begin{itemize}
    \item Retrieve or build the compartment:
    \begin{verbatim}
comp = self.get_compartment(id_compartment)
    \end{verbatim}

    \item Call:
    \begin{verbatim}
temporal.readSimData(...)
    \end{verbatim}

    \item Convert to DataFrame:
    \begin{verbatim}
temporal.simMatrixToDf(...)
    \end{verbatim}

    \item Optionally write a tab-separated CSV.
\end{itemize}

The method returns an \texttt{ExtractValuesResponse} containing:
\begin{itemize}
    \item the DataFrame,
    \item the CSV path (if written),
    \item metadata.
\end{itemize}

\subsection{\texttt{apply\_temporal\_operator(...)}}

\textbf{Purpose:}  
Apply temporal aggregation (annual mean, monthly sum, quantiles, etc.).

It wraps:
\begin{verbatim}
Manage.Temporal.aggregate_matrix(...)
\end{verbatim}

The GUI provides the DataFrame to aggregate, typically produced by \texttt{extract\_values()}.

\subsection{\texttt{compute\_mass\_balance(...)}}

\textbf{Purpose:}  
Compute interannual budgets and mass balances.

Wraps:
\begin{verbatim}
Manage.Budget.calcInteranualBVariable(...)
\end{verbatim}

Returns a \texttt{MassBalanceResponse}.

\subsection{\texttt{apply\_spatial\_operator(...)}}

Provides a unified spatial API.

\begin{itemize}
    \item \texttt{operator = 'catchment\_cells'}  
    Wraps \texttt{Manage.Spatial.getCatchmentCellsIds(...)}  
    Returns a DataFrame with \texttt{id\_cell}.

    \item \texttt{operator = 'aquifer\_outcropping'}  
    Wraps \texttt{Manage.Spatial.buildAqOutcropping(...)}  
    Returns a DataFrame with \texttt{id\_abs}.
\end{itemize}

\section{Persistence and Bancarisation}

HT inherits from \texttt{HTPersistenceMixin}, which provides:
\begin{itemize}
    \item \texttt{HydrologicalTwin.to\_pickle(path)}
    \item \texttt{HydrologicalTwin.from\_pickle(path)}
\end{itemize}

This enables:
\begin{itemize}
    \item caching of fully built hydrological twins,
    \item fast reloads in new QGIS sessions,
    \item the creation of a future \emph{HydroTwin Hub} of \texttt{.pkl} snapshots.
\end{itemize}

Each snapshot stores:
\begin{itemize}
    \item a schema version,
    \item the class name,
    \item the serialized HT instance.
\end{itemize}

\section{Interface–Backend Contract}

The QGIS plugin must interact only with:
\begin{itemize}
    \item \texttt{HydrologicalTwin}
\end{itemize}

and never directly with:
\begin{itemize}
    \item \texttt{Compartment}
    \item \texttt{Mesh}
    \item \texttt{Observations}
    \item \texttt{Manage.*}
\end{itemize}

This guarantees:
\begin{itemize}
    \item a clear API contract,
    \item easier refactoring,
    \item easier automated testing,
    \item reuse by other tools (CLI, notebooks, web services).
\end{itemize}



\section{Current Functional Scope of HydroTwin on CaWaQS Viz}


This section summarizes the functionalities already implemented in CaWaqS viz, that ar e being kept during the work to make this as an hydrotwin.
It reflects the present development state and serves as a reference for ongoing architectural and API design decisions.

\subsection{Simulation vs Observation Analysis}

\begin{itemize}
  \item Comparison for \textbf{aquifer (AQ)} or \textbf{hydrological (HYD)} variables
  \item Distinction between:
  \begin{itemize}
    \item Period used for plotting
    \item Period used for statistical computation
  \end{itemize}
  \item Supported statistics:
  \begin{itemize}
    \item Minimum, maximum, mean, quantiles
  \end{itemize}
  \item Probability density function (PDF) plotting
  \item Interactive graphical outputs
\end{itemize}

\subsection{Spatial Maps}

\begin{itemize}
  \item Aggregation period selection
  \item Temporal frequency:
  \begin{itemize}
    \item Annual, monthly, daily
    \item Multi-annual aggregation for annual and monthly frequencies
  \end{itemize}
  \item Aggregation operators:
  \begin{itemize}
    \item Cumulative, average, quantile, minimum, maximum
  \end{itemize}
  \item Supported variables:
  \begin{itemize}
    \item Rainfall
    \item Effective rainfall
    \item Actual evapotranspiration (ETR)
    \item Infiltration (INF)
    \item Runoff
    \item Overland flow (OZ)
    \item River--aquifer exchanges
    \item Groundwater recharge
    \item Surface overflow
  \end{itemize}
\end{itemize}

\subsection{Water Budget Analysis}

\begin{itemize}
  \item User-defined plotted period
  \item Temporal frequency:
  \begin{itemize}
    \item Annual, monthly, daily
    \item Multi-annual aggregation for annual and monthly frequencies
  \end{itemize}
  \item Aggregation operators:
  \begin{itemize}
    \item Cumulative, average, minimum, maximum
  \end{itemize}
\end{itemize}

\subsection{Hydrological Regime Analysis}

\begin{itemize}
  \item User-defined plotted period
  \item Variables:
  \begin{itemize}
    \item Discharge
    \item Piezometric head
  \end{itemize}
  \item Statistical plots:
  \begin{itemize}
    \item Static outputs (PNG, PDF)
    \item Interactive plots
  \end{itemize}
\end{itemize}

\subsection{Statistical Performance Criteria}

\begin{itemize}
  \item Analysis for:
  \begin{itemize}
    \item Piezometric (PZ)
    \item Hydrological (HYD) variables
  \end{itemize}
  \item User-defined plotting period
  \item Aggregation operators:
  \begin{itemize}
    \item Mean, minimum, maximum, quantiles
  \end{itemize}
  \item Implemented criteria:
  \begin{itemize}
    \item Percent bias (PBIAS)
    \item Average simulated / observed ratio
    \item Kling--Gupta Efficiency (KGE)
    \item Root Mean Square Error (RMSE)
    \item Nash--Sutcliffe Efficiency (NSE)
  \end{itemize}
\end{itemize}


